\chapter{Related Work and State of the Art}
\label{chap:related}

This chapter reviews the three main research areas on which this thesis builds: quantum repeater protocols, entanglement purification, and \Gls{rgs}-based architectures. It then places the scheduling gap identified in Section~\ref{sec:intro:gap} in the context of the existing state of the art.

\section{Quantum Repeater Protocols}
\label{sec:rw:repeaters}

Sending a photonic qubit directly through optical fiber or free space is limited by the channel's transmissivity, i.e. the probability that a photon reaches the receiver. For a channel of length $L$ and attenuation coefficient $\alpha$, this is $\eta=e^{-\alpha L}$. At the typical telecom C-band attenuation of $\alpha_{\mathrm{dB}}\approx0.2~\mathrm{dB/km}$~\cite{ITUT2024Cable}, the transmissivity is already as low as $\eta\sim10^{-40}$ over 2000~km. Thus, even with perfect sources and detectors, direct transmission becomes impractical over long distances, with the achievable rate decreasing exponentially with distance~\cite{Sangouard2011Quantum}. More generally, the \emph{repeaterless secret-key capacity} of a lossy channel, bounded by the \Gls{plob}, has the same exponential scaling in the high-loss regime~\cite{Pirandola2017Fundamental,Pant2017Ratedistance,Azuma2023Quantum}.

A quantum repeater~\cite{Briegel1998Quantum} divides a long channel into shorter elementary links, establishes entanglement over those links, and stitches them into an end-to-end pair through entanglement swapping~\cite{Zukowski1993Eventreadydetectors,Grudka2008Entanglementswapping}. Repeater proposals are commonly grouped into three generations according to how they address photon loss and operational errors~\cite{Muralidharan2016Efficient,Azuma2023Quantum}. First-generation (1G) repeaters use \Gls{heg} and entanglement purification~\cite{Bennett1996Purification,Deutsch1996Quantum} (Section~\ref{sec:rw:purification}), storing generated pairs in matter-qubit memories while other links are established~\cite{Briegel1998Quantum,Duan2001Longdistance}. Second-generation (2G) repeaters replace purification with \Gls{qec} for operational errors while retaining heralding for photon loss~\cite{Jiang2009Quantum,Wo2023Resourceefficient}. Third-generation (3G) repeaters use quantum error-correcting codes~\cite{Knill2000Theory,Gottesman1997Stabilizer} to protect against both loss and operational errors, avoiding heralding across every elementary link~\cite{Fowler2010Surface,Azuma2015All}.

The all-photonic quantum repeaters~\cite{Azuma2015All} studied in this thesis belong to the third generation but differ from proposals based on matter-qubit error correction. They encode the transmitted qubit in a multi-photon graph state~\cite{Briegel2001Persistent,Hein2006Entanglement,Hein2004Multiparty}: the \Gls{rgs}~\cite{Azuma2015All}. Its tree encoding provides loss tolerance while operating on flying qubits, removing the need for quantum memories at intermediate stations. Section~\ref{sec:rw:rgs} places this architecture among related proposals; Chapter~\ref{chap:background} gives the mechanics required here.

A related but architecturally distinct proposal, the measurement-based quantum repeater~\cite{Zwerger2012Measurementbased}, stores small pre-prepared entangled resource states~\cite{Raussendorf2002oneway,Raussendorf2002Computational,Raussendorf2003Measurementbased,Nielsen2006Clusterstate,Greenberger2009Compendium} at each station and couples them to incoming photons using simple \Glspl{bsm}, rather than constructing one graph state spanning the repeater chain~\cite{Browne2005Resourceefficient,Ewert20143,Zwerger2018Longrange}. Its emphasis on minimising the number of qubits held at each station provides a useful comparison for the resource costs studied in Chapter~\ref{chap:results}.

\section{Entanglement Purification Protocols}
\label{sec:rw:purification}

Chapter~\ref{chap:intro} introduced the basic trade-off between copies and fidelity; this section covers the protocols implementing it. The two foundational two-qubit protocols, both using only \gls{locc}, are BBPSSW~\cite{Bennett1996Purification}, which operates on Werner states, and DEJMPS~\cite{Deutsch1996Quantum}, which extends the approach to general Bell-diagonal states without first depolarizing them. Both consume two copies and return one, succeeding only for appropriate local measurement outcomes. Thus, purification improves the quality of an entangled pair by consuming additional entangled resources, directly creating the resource-fidelity trade-off relevant to this thesis.

Entanglement pumping applies the same principle repeatedly: a working pair is purified against a sequence of freshly generated pairs, increasing its fidelity with each round~\cite{Dur1999Quantum,Deutsch1996Quantum,Bennett1996Purification}. In contrast to applying a fixed purification circuit once, pumping allows resources to be accumulated over multiple rounds while retaining one working pair. Both the resource/fidelity trade-off and pumping recur in the \Gls{rgs}-specific circuits on which this thesis' formal model is based (Section~\ref{sec:bg:circuits}).

Later work treats the purification circuit itself as an optimisation problem rather than fixing it in advance. These approaches search over local Clifford operations for shorter circuits with higher success probability and output fidelity than the BBPSSW/DEJMPS family~\cite{Krastanov2019Optimized}, particularly under finite resources and imperfect local operations. Purification has also been extended from bipartite Bell pairs to multipartite graph states~\cite{Dur2003Multiparticle,Kruszynska2006Entanglement}. The \Gls{rgs}-specific two-to-one circuits used by the Integrating paper (Section~\ref{sec:bg:circuits}) adapt these principles to the \Gls{rgs} architecture rather than directly implementing BBPSSW or DEJMPS.

\paragraph{Photon loss tolerance} Purification addresses noise, not photon loss. Tree encoding~\cite{Varnava2006Loss} provides a different form of redundancy: instead of consuming multiple copies to improve fidelity, it encodes one logical qubit in a tree-shaped graph state whose branching factors determine its loss tolerance. Logical measurement outcomes can then be reconstructed from surviving branches. The approach was introduced for one-way measurement-based quantum computation~\cite{Raussendorf2003Measurementbased,Raussendorf2002oneway,Varnava2006Loss} and later developed into a loss-tolerant quantum memory scheme~\cite{Varnava2007Loss}, together with an analysis of its resource-efficiency threshold~\cite{Varnava2008How}. This tree encoding, rather than purification, allows the inner qubits of the \Gls{rgs} to tolerate photon loss (Section~\ref{sec:bg:hrgs}).

\section{Repeater Graph State Architectures}
\label{sec:rw:rgs}

The \Gls{rgs} was proposed as the resource state for the first all-photonic quantum repeater: a graph state whose outer qubits form entangling links to neighbouring stations while its inner qubits are tree-encoded~\cite{Varnava2006Loss}, allowing the logical connection to survive photon loss~\cite{Azuma2015All}. Later work has focused on efficient \Gls{rgs} generation and tolerance to photon loss and operational errors. The original proposal generates the graph state by fusing small entangled photon sets using linear optics~\cite{Azuma2015All,Pant2017Ratedistance}, while deterministic emitter-based schemes build it directly from a small number of photon emitters, reducing the physical device count per station~\cite{Borregaard2020OneWay,Li2022Photonic,Ghanbari2024Optimization}. These approaches have also been demonstrated experimentally at small scale~\cite{Thomas2024Fusion}. For loss and error tolerance, alternative graph topologies and adaptive measurement strategies can improve entanglement rates for a fixed loss and error budget~\cite{Patil2024Improved,Li2025Generalized,Zhan2023Performance}; code-concatenated variants instead address both through two specialized codes rather than a single tree code~\cite{Wo2023Resourceefficient}.

The Bridging and Integrating papers, discussed in Section~\ref{sec:intro:gap} and formalised in Chapter~\ref{chap:background}, provide the starting point for this thesis. The former introduces the \Gls{hrgs} building block~\cite{Benchasattabuse2026Bridging,Benchasattabuse2026Resource,Benchasattabuse2024Architecture}, while the latter shows that purification can be incorporated using optimistic (blind) execution~\cite{Mobayenjarihani2023Optimistic,Benchasattabuse2025Integrating}. This establishes that purification can be integrated, but not how it should be scheduled: where and how often to purify across a full chain remains open, as discussed in Section~\ref{sec:rw:positioning}.

\section{Schedule Optimisation in Quantum Networks}
\label{sec:rw:scheduling}

Deciding when and where to apply entanglement purification is a separate problem, even when purification is available throughout a repeater chain. In adaptive schemes, a repeater node selects its next operation based on fidelity or heralding outcomes observed so far rather than following a fixed plan~\cite{Munro2015Quantum}. Reinforcement learning and other data-driven controllers have also been proposed for quantum-network control. One study, for example, reports substantial offline costs in richer settings, including exponentially growing state representations, training times of up to $\sim 36$ hours per configuration, and hyperparameter searches needed to outperform a heuristic baseline~\cite{Yau2026Reinforcement}.

Optimistic purification has likewise been studied in generic quantum networks, independently of any \Gls{rgs}-specific construction~\cite{Mobayenjarihani2023Optimistic}, showing that deferred heralding can be useful beyond the architecture studied here. Another scheduling decision is how qubits of a logical channel are multiplexed across serial or parallel physical links, which affects end-to-end rate and fidelity~\cite{Meter2007Communication}. These results demonstrate that the ordering and placement of operations can influence network performance, even when the underlying physical operations are already fixed.

\section{Positioning of This Work}
\label{sec:rw:positioning}

The work surveyed here generally optimises a fixed repeater component, such as a purification circuit or graph-state construction, or evaluates a specified scheduling policy. For the \Gls{hrgs} architecture studied in this thesis, however, these decisions interact: the choice of purification location affects the resources consumed, the resulting fidelity, and the timing of subsequent operations. No systematic search considers where to purify, how many copies to consume, which circuit to use, and in what order across a full chain. As Section~\ref{sec:intro:gap} establishes, this is the gap addressed by this thesis.

\TODO{since shorter, include image of RGS? and/or other figures.... to design myself}